
\documentclass[12pt]{article}
\usepackage{amsmath, amssymb, geometry}
\usepackage{booktabs}
\geometry{margin=1in}
\title{\textbf{SAT.4D Supplementary Definitions, Falsifiability, and Interpretive Modes}}
\author{SAT Framework Team}
\date{June 2025}

\begin{document}
\maketitle

\section*{1. Core Structural Quantities}

\subsection*{1.1 Misalignment Angle $\theta_4(x)$}
Defined as the local angular misalignment between a filament’s tangent vector $v^\mu$ and the foliation normal $u^\mu$:
\[
\cos(\theta_4(x)) = \frac{v^\mu u_\mu}{\|v^\mu\| \, \|u^\mu\|}
\]
\textbf{Interpretation:} Measures resistance to null propagation across $\Sigma_t$. Emerges during wavefront intersection and reflects local inertial signature.

\subsection*{1.2 Strain Tensor $S_{\mu\nu}(x)$}
Defined via the gradient of the foliation vector field:
\[
S_{\mu\nu}(x) = \nabla_\mu u_\nu + \nabla_\nu u_\mu
\]
\textbf{Interpretation:} Encodes local foliation distortion. Acts as the geometric source of curvature in emergent GR.

\subsection*{1.3 Twist Sector $\tau(x)$}
A scalar or multi-component field capturing the rotational embedding of filaments in local 3D slices. Represents helicity or internal quantum number structure.

\textbf{Interpretation:} Governs spin-class behavior, chirality, and inter-filament coupling dynamics.

\section*{2. Interpretive Modes of SAT.4D}

\begin{center}
\begin{tabular}{@{}llll@{}}
\toprule
\textbf{Mode} & \textbf{Name} & \textbf{Block Dynamics} & \textbf{Use Case} \\
\midrule
1 & True Block & None (fully static 4D) & Core derivations, structure-only \\
2 & Semi-Dynamic & Local adjustments at $\Sigma_t$ & Quantum interpretation, entropy \\
3 & Solidification Front & Structure resolves at $\Sigma_t$ & Cosmology, time asymmetry \\
\bottomrule
\end{tabular}
\end{center}

\textit{Note:} All SAT module construction must occur strictly in Mode 1.

\section*{3. Falsifiability Table}

\begin{center}
\begin{tabular}{@{}lll@{}}
\toprule
\textbf{Prediction} & \textbf{Testable Condition} & \textbf{Consequence if Violated} \\
\midrule
$Q \leq 3$ bound states only & Observation of stable $n \geq 4$ particles & SAT.O4 falsified \\
Mass $\propto 1/Q$ & Deviations from $m_i/m_j \sim Q_j/Q_i$ & Topological suppression invalid \\
Quantized $\theta_4$, $S_{\mu\nu}$ & No geometric signature of mass & Curvature emergence falsified \\
No SUSY partners & Detection of sleptons/gauginos & SAT particle ontology refuted \\
Gauge couplings from $\rho$ & Inconsistent coupling ratios & SAT.O5 falsified \\
\bottomrule
\end{tabular}
\end{center}

\section*{4. Topological Quantum Field Theory in SAT.4D}

\subsection*{4.1 Topological Quantization}
Quantization arises from discrete topological invariants:
\begin{itemize}
  \item Winding number $n$
  \item Linking number $L_k$
  \item Knot complexity (e.g., Jones polynomials)
\end{itemize}

\subsection*{4.2 Wavefunctions and Coherence}
Define phase function:
\[
\Psi[\text{Topology}] \sim \exp(i\theta[\text{Linking, Winding, Knotting}])
\]
Analogous to a quantum wavefunction but structurally grounded in filament configurations.

\subsection*{4.3 Entanglement and Superposition}
Entanglement: Nonlocal linking of filament bundles.\\
Superposition: Overlapping topological states with shared macroscopic projection.

\subsection*{4.4 Emergent Path Integral}
\[
Z = \int \mathcal{D}\gamma \, e^{-T \int \|\ddot{\gamma}\|^2 d\lambda}
\]
Quantifies structural amplitude across topological bundle space.

\subsection*{4.5 Observables}
All quantum observables (e.g., energy, spin, charge) emerge from:
\[
\mathcal{O}[\gamma] \in \{Q, \theta_4, \tau, S_{\mu\nu}\}
\]

\end{document}
